% =========================================================
% Proof Stub: Convergence Criterion via Limsup and Liminf
% Source: volume-iii/analysis/sequences/notes/asymptotics/notes-limsup-liminf.tex
% Mode: standard
% =========================================================

\clearpage
\phantomsection
\hypertarget{proof-RA-ABB-T-limsup-conv-criterion}{}

\subsubsection[Convergence Criterion via Limsup and Liminf]{Proof --- Convergence Criterion via Limsup and Liminf}

\bigskip

\noindent
\textbf{Source.}
Abbott, \emph{Understanding Analysis}; see \texttt{notes-limsup-liminf.tex}.

\vspace{0.75em}

\noindent
\textbf{Goal.}
$(x_n)$ converges iff $\limsup x_n = \liminf x_n$, and the common value is the limit.

\vspace{0.75em}

\noindent
\textbf{Logical form.}
\[
  \lim x_n \text{ exists} \Longleftrightarrow \limsup x_n = \liminf x_n.
\]

\vspace{0.75em}

\noindent
\textbf{Proof strategy.}
(⇒) If $x_n \to L$, show both $s_n$ and $i_n$ converge to $L$ using the $\varepsilon$-$N$ definition. (⇐) If $\limsup = \liminf = L$, squeeze: $i_n \le x_n \le s_n$ and both ends go to $L$.

\vspace{0.75em}

\noindent
\textbf{Proof.}
\begin{proof}
\Claim $(x_n)$ converges iff $\limsup x_n = \liminf x_n$, and the common value is the limit.

\Given 

\Goal 

\Strategy (⇒) If $x_n \to L$, show both $s_n$ and $i_n$ converge to $L$ using the $\varepsilon$-$N$ definition. (⇐) If $\limsup = \liminf = L$, squeeze: $i_n \le x_n \le s_n$ and both ends go to $L$.

\medskip

% --- (⇒) given ε: for n ≥ N, x_k ∈ (L-ε, L+ε) for all k ≥ n ---
% --- so s_n = sup_{k≥n} x_k ≤ L + ε and i_n ≥ L - ε ---
% --- both s_n → L and i_n → L ---
% --- (⇐) i_n ≤ x_n ≤ s_n; squeeze theorem: x_n → L ---

\end{proof}

\begin{remark*}[Proof shape]
(⇒) $\varepsilon$-$N$ forces tail sup and inf to converge to the same limit. (⇐) Squeeze between $i_n$ and $s_n$.
\end{remark*}

\begin{remark*}[Dependencies]
\begin{itemize}
  \item Squeeze Theorem.
  \item $\varepsilon$-$N$ definition of convergence.
  \item $i_n \le x_n \le s_n$ for all $n$ (definition of inf and sup of tail).
\end{itemize}
\end{remark*}
